\chapter{Detailed Design and Implementation}
\pagestyle{fancy}

\subsection{Users and Use Cases}

From the Vision we can derive that we have two types of users. The main one will simply be called \b{User}. By this, we refer to people who download the software, set it up and want to use it. One such user can use the system for many use cases, which can be grouped into categories:

\begin{itemize}
    \item Chess piece image acquisition
    
    We need to get images of chess pieces. To achieve this, a previous configuration will be necessary, to get the corners of the board (\i{Configure}), after which these cell images can be extracted and labeled (\i{Crop and Label}). To be more user-friendly, test variations of these functions shall be provided as well, such that the results of a potential Configure or Crop action can be seen without any effects on previous values.
    
    \item Chess piece image classification

    After having the piece images, we can use Machine Learning algorithms to classify them as either White Free, Pawn, Bishop, Knight, Rook, Queen, King or one of the black counterparts. For this, multiple types of classifiers will be defined, from which the user can choose (\i{Select Classifier}). These can be saved or loaded to save time, trained if necessary and tested to check for metrics. The main functionality \i{Classify Board} can then take a set of 8x8 cell images and classify all of them, reconstructing a board.

    \item Move validation
    
    If we have a board position, we can compare it to the previous one, and see if a valid move happened ``in between''. This is done by \i{Validate Move}. Some moves might need to be discarded, as well as there needs to be an easy way of starting a new game.

    \item Miscellaneous actions
    
    Aside from what has already been listed, the User can also \i{Shuffle and Split} to mix up the cell images, \i{Clear All} to delete all images, or \i{Change Settings} to modify parameters which persist cross-session.
\end{itemize}

The secondary actor of the system is \b{Developer}. This could be any open-source connoisseur willing to contribute to the codebase or reuse parts of it. For this, it is essential to have readable code, good documentation and clear version control.

These two actors and the associated use cases can be seen on Figure \ref{fig:usecase}.


\begin{figure}[H]
    \centering
    \includegraphics[width=0.6\textwidth]{figs/usecase.png}
    \caption{Use Case Diagram}
    \label{fig:usecase}
\end{figure}



% ------------------------------------------------------------------------------------------- %
% --------------------------------- User Application Module --------------------------------- %
% ------------------------------------------------------------------------------------------- %

\section{Modules and Communication}

\begin{figure}[H]
    \centering
    \includegraphics[width=0.6\textwidth]{figs/black-box.png}
    \caption{Black box of the system and its modules}
    \label{fig:black-box}
\end{figure}

\begin{figure}[H]
    \centering
    \includegraphics[width=\textwidth]{figs/signals-slots.png}
    \caption{Important signals and slots}
    \label{fig:signals-slots}
\end{figure}

\section{User Application Module}

The User Application Module provides a Graphical User Interface for the user, as well as serves as the main module, given that our application is a Qt-based desktop application. As a result of this, class \t{UserApplication} is more than just a wrapper of other UI-related classes: the main loop of the whole application is the event loop provided by its \t{run} method. The UI package on itself is simple, see Figure \ref{fig:class-ua}. Two separate windows are available, the main window and the settings window, which will be described in the next paragraphs.

\begin{figure}[H]
    \centering
    \includegraphics[width=0.8\textwidth]{figs/class_ua.png}
    \caption{Classes of package \t{ua}}
    \label{fig:class-ua}
\end{figure}

As mentioned before, everything starts in \t{UserApplication}. This application uses an instance of the \t{MainWindow} class. Figure \ref{fig:main-window} shows how this appears to the user. \t{MainWindow} inherits from \t{QMainWindow}, and is a \t{Q\_OBJECT}, since it has signals and slots. The goal is to have a simple UI, so I tried to use the simplest and most suitable widgets to build up this window. Unmodifiable texts (text which might change through code but not through the user directly editing it) are displayed in \t{QLabel}s. Buttons are of type \t{QPushButton}. Dropdowns in Qt come in the form of \t{QComboBox}es. And with some trickery, these elements were enough.

\begin{figure}[h]
    \centering
    \includegraphics[width=\textwidth]{figs/main-window.PNG}
    \caption{The main window}
    \label{fig:main-window}
\end{figure}

All the elements are grouped in layouts so that they are placed as intended. There are three kinds of layouts used, \t{QHBoxLayout} for horizontal, \t{QVBoxLayout} for vertical alignment of widgets, and \t{QGridLayout} for matrix-like organization. These are not declared in the header file. The header contains only the elements that need to be accessed later, and layouts are not part of this category. To make layouts have backgrounds (which is important for aesthetics, but also for debugging), a typical Qt pattern is used, associating each layout with a ``dummy'' widget. The steps for adding widgets to the application/ building up one block of the window typically are:

\begin{enumerate}
    \item Find an existing layout and associated widget, or define new ones. \\
    \t{newWidget = new QWidget();} \\
    \t{newLayout = new QVBoxLayout(newWidget);}

    \item If you chose the second option and it is needed, define the layout's size and color. \\
    \t{newWidget->setFixedHeight(height);} \\
    \t{newWidget->setFixedWidth(width);} \\
    \t{newWidget->setFixedSize(width, height);} \\
    \t{newWidget->setStyleSheet(styleSheetName);}
    
    \item Instantiate your widgets. Different types of widgets have different signature constructors. \\
    \t{newLabel = new QLabel("text");} \\
    \t{newComboBox = new QComboBox();} \\
    \t{newComboBox->addItem("option");} \\
    \t{newButton = new QPushButton("text");}
    
    \item Add the widgets to the layout. \\    
    \t{newLayout->addWidget(newLabel);} \\
    \t{newLayout->addWidget(newComboBox);} \\
    \t{newLayout->addWidget(newButton);} \\

    \item If a new layout was defined, add it to a higher-level layout.\\
    \t{higherLevelLayout->addWidget(newWidget);}
\end{enumerate}

In the project, there are examples of all layouts and widgets necessary. Figure \ref{fig:main-window-skeleton} shows a ``skeleton'' version of the main window - this is the programmer's perspective. Widgets are grouped by the functionality they serve to make the GUI more intuitive.

\begin{figure}[h]
    \centering
    \includegraphics[width=\textwidth]{figs/main-window-skeleton.png}
    \caption{The skeleton of the main window}
    \label{fig:main-window-skeleton}
\end{figure}

The camera previews are shown in \t{cameraOneImageLabel} and \t{cameraTwoImageLabel}. As their name suggests, they are \t{QLabel}s, but without text. One may set a picture as ``the background'' of a label by calling \t{label->setPixmap(pixmap)}, where pixmap is of type \t{QPixmap}, which we can get from a \t{QImage}. The flow of updating the image labels is the following:
\begin{enumerate}
    \item \t{previewImageReadySignal(QImage(...))} in \t{CameraHandler (IP)}
    \item 
\end{enumerate}

% ------------------------------------------------------------------------------------------- %
% --------------------------------- Image Processing Module --------------------------------- %
% ------------------------------------------------------------------------------------------- %

\section{Image Processing Module}

In the source code, the Image Processing Module can be found in a package/folder called \t{ip} that is responsible for everything related to images: getting, processing, saving and deleting them. The main class \t{ImageProcessing} provides an entry point for all these functionalities. The \t{ImageProcessing} does not do lots of computations, it rather serves as a wrapper for other, more specific IP classes. So the main goal is to wait for user requests in the form of signals from UA and call the appropriate subcomponent. 

The main classes of IP can be seen in Figure \ref{fig:class-ip}. Here, the simple arrows show the ``uses'' relation, for example, \t{ImageProcessing} uses the \t{Configurer} class, \t{Configurer} uses the \t{BFS} class. The empty-headed arrows show ``implements'' relations, like \t{CxxClassifier} implements \t{AbstractClassifier}, \t{KNearestNeighbors} implements \t{CxxClassifier}.

\begin{figure}[H]
    \centering
    \includegraphics[width=\textwidth]{figs/class_ip.png}
    \caption{Classes of package \t{ip}}
    \label{fig:class-ip}
\end{figure}

Now let's take a more detailed look at each subcomponent and see what functionalities they provide.


% --- Camera Handling --- %

\subsection{Camera Handling}

Camera handling is the process of dealing with cameras with the help of OpenCV, listening to requests for images, fulfilling these requests, and continuously providing images for the GUI as a preview of what cameras see.

An instance of the \t{CameraHandler} is capable of handling one camera. Its field index is directly used as an input for OpenCV's \t{VideoCapture} object. The field \t{isRunning} is self-explanatory. Field \t{isImageRequested} is the flag through which a request to save a frame can be made. After constructing, the \t{doWork} function shall be called. It is expected that a \t{CameraHandler} instance is run in a separate thread, which executes this \t{doWork} function. 

The steps of getting an image through \t{CameraHandler} are the following:

\begin{itemize}
    \item Create a \t{SignalWaiter} instance: pass the \t{CameraHandler}, where to save the image, and additionally a \t{signalWaiterName}, for debugging.
    \item Call \t{start} on the \t{SignalWaiter} instance.
    \item Check the return value of start, if 0, continue.
    \item Try reading the image from the path given by the \t{SignalWaiter} instance. If successful, done.
\end{itemize}

For convenience, all this is wrapped in \t{ImageProcessing}'s \t{getImage} method. You may ask, what is all this hassle for? Well, there are many ways in which getting an image can go wrong. The request may be made and a response might never come, or come too late. For this, there is a timer embedded in \t{SignalWaiter}. Also, getting the image is the same for every process, but what we do with it, including where we save it, differs in every use case. So this is a unified way of getting the image, which can then be moved or processed in any way. The big goal is to reduce the overhead of \t{CameraHandler}, since it also has to send frames periodically to the GUI, and low frame rates are visible to the human eye and very annoying.

Since concurrency might appear, access to \t{isImageRequested} is synchronized with \t{isImageRequestedLock}. Signal \t{previewImageReadySignal} is emitted periodically, without waiting for a request. If we have a frame, we send it to the preview panel. There will be two \t{CameraHandler} instances, one for each camera. While requests might come somewhat at the same time, the preview frames do not need to be synchronized.


% --- Configuration --- %

\subsection{Configuration}

Configuration refers to getting the corners of the chessboard. This is a preliminary step since to get the images of cells, we first need to know the corner positions and the width of the chessboard border. Since other phases of image processing highly depend on this step, it is better to choose a less fancy but working solution rather than force ideas that sometimes might work, and sometimes not.

Class \t{Configurer} represents a subcomponent which wraps every functionality through three functions:
\begin{itemize}
	\item \t{configure}: gets the corners of the chessboard
    \item \t{cropAndLabel}: cuts the board image into cells and puts them in folders based on provided information (known encodings, which will come from the user, as later explained)
    \item \t{prepareCellImages}: cuts the board, but this time to prepare for the classification of the board (unknown encodings, these will need to be guessed)
\end{itemize}

In the following, we will talk about the classes wrapped inside \t{Configurer}, by going through these three functions and what happens in them.


\subsubsection{Function \t{configure}}
\paragraph{Filtering}

First, we get the source image. This is a color image, however, for Canny, we will need a preprocessed binary image. The first step \t{Configurer} does in \t{configure} is to resize the image to a fixed size (500x500), convert it to grayscale and apply a Gaussian filter to reduce the noise (\t{Filter}). The result is a slightly blurred image as shown in Figure \ref{fig:conf_initial}:

\begin{figure}[ht]
    \centering
    \begin{subfigure}[b]{0.475\textwidth}
        \centering
        \includegraphics[width=\textwidth]{figs/conf1.png}
        \caption{Grayscale image} 
    \end{subfigure}
    \hfill
    \begin{subfigure}[b]{0.475\textwidth}  
        \centering 
        \includegraphics[width=\textwidth]{figs/conf2.png}
        \caption {Filtered image}
    \end{subfigure}
    \caption{Initial steps of configuration}
    \label{fig:conf_initial}
\end{figure}


\paragraph{Conversion, Morphological Operations}

Next, we take the image from the previous step and convert it to a binary image (class \t{Conversion}). While a binary image could be stored in a boolean matrix, or a bitmap, or whatever fanciness, it is easier to use a \t{Mat\_<char>}, just as in the case of the grayscale images, but with values ``restricted'' to either 0 (object pixel) or 255 (non-object, background pixel). This way it can be easily visualized with \t{imshow}.

To remove possible background noise, an interesting enhancement was added at a late stage of the project. The notion of image edge pixels is defined as pixels that are at IMAGE\_EDGE\_SIZE (3 pixels) distance from the edge of the image. It is a precondition for the algorithm to have the whole chessboard visible, and to have it on a light uniform background. However, the background might not cover the whole image, and in this case, the BFS step might remove the possible noises: the image edge pixels are iterated, and where a black pixel is found, it is assumed that it is noise, and a BFS is started from it, transforming all the nearby black pixels into white ones.

After this step, to further reduce the noise, we apply the closing operation (class \t{Morphological}). This is no more than a dilation followed by an erosion. This helps to get rid of smaller noise created by illumination and smaller details of the chessboard border which are irrelevant. The results are shown in Figure \ref{fig:prep_canny}.
 
\begin{figure}[H]
    \centering
    \begin{subfigure}[b]{0.475\textwidth}
        \centering
        \includegraphics[width=\textwidth]{figs/conf3.png}
        \caption{Binary image} 
    \end{subfigure}
    \hfill
    \begin{subfigure}[b]{0.475\textwidth}  
        \centering 
        \includegraphics[width=\textwidth]{figs/conf4.png}
        \caption {Image after closing}
    \end{subfigure}
    \caption{Preparing for Canny}
    \label{fig:prep_canny}
\end{figure}

\paragraph{Canny}

Class \t{Canny} contains the implementation of Canny Edge Detection, based on Image Processing, Laboratory Works \cite{ip_labs}. The input of the method is assumed to be an already noise-filtered grayscale image, just as described in previous steps. The magnitude and gradient of pixels from the gradients are extracted, which are achieved by applying convolution with a Sobel matrix, in both x and y directions. The angles are quantized into eight numbers, which later on are used for eliminating pixels that do not represent a local maximum based on their magnitude and neighbors in the gradient and opposite gradient direction.

After this, two-level thresholding is done (on the image of magnitudes), with hard-coded values 100 and 200. These could be adaptive or user-defined, but for now, they are hardcoded. Based on these thresholds, pixels are re-evaluated and given a ``naming'', as shown in Table \ref{tab:two-level-thresholding}. The new value could have been arbitrary, but the chosen convention must be respected further on. These values are nice since if we would like to, we could display the thresholded image, in a way that makes sense.

\begin{table}[H]
    \centering
    \begin{tabular}{|l|l|l|}
    \hline
    \textbf{Range} & \textbf{New Value} & \textbf{Naming} \\ \hline
    {[}0, 100)     & 0                  & non-edge        \\ \hline
    {[}100, 200)   & 50                 & weak-edge       \\ \hline
    {[}200, 255{]} & 255                & strong-edge     \\ \hline
    \end{tabular}
    \caption{Two level thresholding}
    \label{tab:two-level-thresholding}
\end{table}

All this is necessary for the edge extension step, in which weak edge pixels are ``appended'' to the strong edges. This is because if we only take strong edges, we might have disconnected borders, and if we take both weak and strong edges (basically lower the threshold), we might get unwanted pixels in the result. To achieve this, we iterate the pixels of the image, and whenever a new strong edge point is found, we start a new Breadth First Search, which accepts both weak and strong edges, and sets the met weak edges to strong ones. To avoid the big overhead of repeating computations, a \t{visited} matrix is defined, which shows if we already checked a given pixel. After this, we can get rid of the remaining weak edges, as they are (hopefully) irrelevant to us.

\begin{figure}[H]
    \centering
    \includegraphics[width=0.45\textwidth]{figs/conf5.png}
    \caption{Step: Canny Edge Detection}
    \label{fig:canny}
\end{figure}


\paragraph{Hough}

As simple as it might seem, the implementation of the Hough line detection algorithm came with more trouble than expected. First, the signature of the base method was hard to define. After multiple changes, the current form of the method returns a vector of \t{lineRhoTheta} structures. One such structure encapsulates the two integer values describing a line (as explained in Subsection \ref{analysis:hough}), as well as two methods for easy handling of the lines: \t{drawLine} to illustrate the result, and \t{getIntersection} since the intersections of the lines are important for us.

You may ask, why are these intersections important? Well, initially I wanted to use the Hough algorithm to detect all the lines of a chessboard, and the intersection points of these lines would have served as the corner points of the cells. In practice, perfect detection of 18 lines cannot be consistently achieved, especially with my setup. Chess boards also come in all forms and shapes, and most of them have borders (the part around the cells that are not cells, just ``decoration''), so this adds at least 4 unwanted lines in the formulae.

One assumption we can safely make is that the chessboard is square. And that it is the biggest square in the image. So, instead of looking for the lines delimiting the cells, I chose to search for the 20 most dominant lines Hough can find, compute all the intersections of all lines, and search for the points which are in the image, and closest to the corner of the image. This way we get the corner points of the chessboard. The fact that chessboard borders vary in size and often are confusing for Image Processing algorithms led me to the decision to give up finding the thickness of the border, and instead make it the user's job to fine-tune the border width parameters. This is not very elegant, but it works.

\begin{figure}[H]
    \centering
    \includegraphics[width=0.45\textwidth]{figs/conf6.png}
    \caption{Step: Hough Line Detection}
    \label{fig:hough}
\end{figure}


\paragraph{Corner extraction}

The last steps of configuration are done directly in method \t{configure}. We build up a vector of \t{Point2i} intersections, by calling \t{getIntersection} on each permutation of two lines. We discard intersections outside the image. Later on, these points are iterated, to get the points that are closest to the corners of the image. From the size of the image we know these corners, and the distance is simple to get with the Euclidean distance equation. The corner points are stored since we need them in further operations.

\begin{figure}[H]
    \centering
    \begin{subfigure}[b]{0.475\textwidth}
        \centering
        \includegraphics[width=\textwidth]{figs/conf7.png}
        \caption{Getting all intersections} 
    \end{subfigure}
    \hfill
    \begin{subfigure}[b]{0.475\textwidth}  
        \centering 
        \includegraphics[width=\textwidth]{figs/conf8.png}
        \caption {Getting corners}
    \end{subfigure}
    \caption{Intersection of lines from Hough}
    \label{fig:corner-extraction}
\end{figure}


\subsubsection{Function \t{cropAndLabel}}

This function has the purpose of building a training set. The user sets up the chess GUI to resemble what he/ she sees on the board, then presses the \t{cropAndLabel} button. The information from the chess GUI is extracted and passed on to this function through the \t{encodings} parameter. The function has to crop \t{imgOriginal} into cell images and give them names. If parameter \t{isTest} is false, a save operation request is made for each cell image, with the correct folder path. Otherwise, the cell images are not saved, they are just displayed, enabling the user to test if border sizes are appropriate without messing up the dataset. These functionalities map to two different buttons on the GUI.

At this point, we must assume that corner extraction worked well, and we have the four corner points available. By calling a utility function \t{warpAndRemoveBorder}, we extract from the image the warped region of interest. This is done by using OpenCV's \t{getPerspectiveTransform} and \t{warpPerspective} functions. From the resulting image, we can remove the borders, since the border dimensions are given by the user (or default values are used). The results look like in Figure \ref{fig:rectification}.

\begin{figure}[H]
    \centering
    \begin{subfigure}[b]{0.475\textwidth}   
        \centering 
        \includegraphics[width=\textwidth]{figs/conf9.png}
        \caption{Warping the board}
    \end{subfigure}
    \hfill
    \begin{subfigure}[b]{0.475\textwidth}   
        \centering 
        \includegraphics[width=\textwidth]{figs/conf10.png}
        \caption{Removing the border}  
    \end{subfigure}
    \caption{Making use of corner points and border dimensions}

    \label{fig:rectification}
\end{figure}

The name of the files could be arbitrary, as long as it is not repeating (otherwise images would overwrite each other). I chose to compose a meaningful name, so it could help with debugging. The name contains the row and column of the cell, in the main coordinate system, alongside a timestamp to ensure uniqueness. This is achieved through mapping the camera's coordinate system, as explained in Representation and Coordinate Systems \ref{analysis:representation}. The folder path comes from the parameter \t{encodings}, and \t{FileHandler}'s \t{saveImage} utility is called, which will be described later. Some examples of cell pictures can be seen in Figure \ref{fig:cells}.

\begin{figure}[H]
    \centering
    \includegraphics[width=0.45\textwidth]{figs/conf11.png}
    \caption{Cell image examples}
    \label{fig:cells}
\end{figure}


\subsubsection{Function \t{prepareCellImages}}

This function is very similar to the previous one since it also crops the cell images and saves them. However, this time, the images are saved to a different path, defined by \t{Paths::BOARD\_FOLDER}, and in their name only contains the string ``board'' and the row and column. This convention is used so that the classifiers can read them later and know which image is tied to which cell; they can later unite the classifications and return a matrix of predictions.


% --------------------------------- Classification --------------------------------- %

\subsection{Classification}

Classification refers to the classification of cell images and preceding operations, like training and testing the active classifier. This subcomponent's wrapper is called \t{Classifier}. It holds one instance of each classifier: \t{KNearestNeighbors}, \t{SupportVectorMachine} and \t{ConvolutionalNeuralNetwork}. These classifiers implement the \t{AbstarctClassifier} class. Since there were multiple similarities between KNN and SVM implementation, both being written in C++ and directly using OpenCV, the intermediary \t{CxxClassifier} was added to avoid repeating code.

At each moment, one of the three classifiers is selected as active. By default, this is KNN. If changes are made to the classification parameters, the affected classifiers are re-instantiated. The user might make a backup for these classifiers to avoid re-training them each session, however, currently, only one backup per type is allowed.


% --- K Nearest Neighbors --- %

\subsubsection{K Nearest Neighbors}

This classifier is the simplest option a user can choose. It is based on the KNN presented in Pattern Recognition Systems. For features, it uses the color histogram, more precisely the three color histograms (blue, green, red) with a given resolution, put ``one after another''. The resolution of the histograms is given by the \t{numberOfBins} parameter, which by default is set to 8. This means that the color value of one channel, ranging from 0 to 255, gets mapped into one of the eight bins, based on where it lies on the scale. This will result in a feature of length $8 x 3 = 24$.

Such a feature gives acceptable results if the colors of the image hold valuable information regardless of the underlying shape. In our case, we can assume that images of pieces of the same class will have similar features, but we cannot expect a really good outcome, since with such features we completely use the notion of dimensionality. One can imagine that a White Free and a Black Free cell would result in very similar outcomes since one is the vertically flipped version of the other. However, this particular mistake is not a mistake: in the long run, we do not care about the color of a free cell, since this can be deduced from its position, and as the rules of chess go, a free cell is a free cell regardless of its color.

Then why collect white and black free cells in differently labeled folders? Well, for other classification algorithms, this might be desirable. In the testing section, the metrics and outcomes will be discussed. We will see how, to make a fair comparison, at the moment of evaluation white and black free cells will be considered the same class, even if during classification they were separated. The parameter \t{uniteFrees} lets the user decide whether to put these into one category from the beginning or only after the classification, but in the final ``merged'' results it does not make much difference.

Depending on this \t{uniteFrees} parameter, the internal labels used by KNN might lie in the range [0,12] or [0,13]. Mapping happens at an upper level, at \t{CxxClassifier}'s \t{externalToInternal} and \t{internalToExternal} methods; this code will be reused in the SVM source code.

During training, the images are loaded, transformed into features, and stored. This is not a conventional training, KNN can be considered a ``lazy'' model, and a very memory-hungry one. While this step might go relatively fast, at the moment of prediction, a distance from each input image's feature has to be calculated and sorted. With the increase in the number of images, the waiting time increases considerably as well. 

Overall, I chose this model to have one classifier done fast, to help me think of how it needs to be integrated and which are the functionalities that could be extracted into an interface, to deal with changes (i.e. addition of new models, low coupling). The KNN algorithm is not necessarily the one to get blamed here, the choice of feature is what makes this solution inferior compared to SVM and CNN.


% --- Support Vector Machine --- %

\subsubsection{Support Vector Machine}

While it looks like support vector machines are getting out of fashion, they still hold their ground. Implementing an SVM involves high-level mathematics, which comes with the risk of introducing hard-to-find mistakes and requires extensive study. Instead of implementing my own SVM and writing code to extract features to feed it, I chose to experiment with OpenCV's implementation of SVM and the ORB features.

Unfortunately, OpenCV's features and classifiers are poorly documented, and it is hard to find resources that truly explain how to use them. A kind Stack Overflow user explains \href{https://stackoverflow.com/questions/29870726/brief-implementation-with-opencv-2-4-10}{in a question thread} that using one such classifier usually involves the following steps:

\begin{enumerate}
    \item \b{Finding some interesting (feature) points in an image.} Usually, this is done with \t{FeatureDetector}s that return a vector of key points based on different criteria. I experimented with SIFT detectors using the \t{SIFT::create} method. The problem was that this yielded different amounts of key points for different images. This means variable length features, which is not desirable. The workaround suggested by the assistant teacher was to hardcode a grid of key points. While this might sound weird at first, it makes sense: we will look at fixed-distance pixels and see what ``significance'' they have.
     
    \item \b{Finding a way to describe the key points.} For this task, \t{DescriptorExtractor}s come in serve. Trial and error led to the choice of ORB. Such a descriptor can be instantiated using \t{ORB::create}, returning a \t{Ptr<ORB>}, which is OpenCV's implementation of smart pointers. We can then directly call \t{compute} on this instance, giving us a descriptor (we previously called this a feature vector).
    
    \item \b{Using the features and feeding it to an algorithm.} While OP mentions feature matching using \t{DescriptorMatcher}, I found a way to feed these feature vectors to the OpenCV SVM. We get an SVM instance with \t{ml::SVM::create}, and set different parameters, like a linear kernel, a maximum number of iterations and an accuracy threshold.
\end{enumerate}

To compile and run all this, the \t{opencv2/ml.hpp} header had to be included, alonside the \t{opencv2/xfeatures2d.hpp} header, which required \t{opencv4[contrib]} installation. There were some sacrifices I had to make to fit all the premade interfaces. For example, the features had to be of type float for SVM to accept it. This meant that I could not fit the feature matrix and class labels into \t{CxxClassifier}, and I had to write some very similar code in both the KNN and SVM source code due to slight type differences. While the used SVM has many high-level functions, for evaluation I chose to go with a naive implementation of calling predict separately on each image and building my confusion matrix, since the available functions were too high-level and not transparent enough.


% --- Support Vector Machine --- %

\subsubsection{Convolutional Neural Network}

Although there are C++ alternatives, for implementing a Convolution Neural Network I chose to use Python and Keras/Tensorflow. This means a slight overhead since the C++ code has to call a Python script (through a \t{system} call). While in some cases I found it enough to use the console for printing messages, in the case of board classification and testing it was necessary to return more than just an integer to the C++ code. This was done through files, which is another overhead, but since we talk about small files and C++, this mostly goes unseen.

The three scripts \t{cnn\_train.py}, \t{cnn\_test.py} and \t{cnn\_classify\_board.py} along with \t{cnn\_common.py} for methods and libraries used in multiple scripts compose the CNN classifier. To successfully run these scripts, one needs to set up an environment with the proper packages and versions, which often is a hassle. GPU usage is recommended as well, although we have a small model here.

\begin{enumerate}
    \item \b{cnn\_train.py} takes five arguments: the training images' directory, the validation images' directory, the directory where the model needs to be saved, the number of epochs training should go on and if augmentation shall be applied. The last two are user-defined parameters. The augmentation layers include random horizontal flipping (vertical makes no sense in this context), slight random rotations, and contrast and brightness changes. Since we are in the Python world, we can use \t{pyplot} to visualize the data and some sample augmentations. The model is taken directly from the Tensorflow tutorials page. From my experience, changing and fine-tuning such models requires advanced knowledge of CNNs and plenty of time, usually with not-so-satisfying results from amateurs like me. The model is saved implicitly since the other scripts need to load it somehow.
    
    \item \b{cnn\_test.py} loads the images in a similar fashion as \t{cnn\_train.py}. Evaluation of the model is done with \t{predict}, however, for building the confusion matrix, a \t{predict} call was needed as well. With predictions and real labels one near another, the confusion matrix was manually built to obey the rules imposed by \t{AbstractClassifier}. Special care had to be taken since the internal encodings were based on the folder names ordered alphabetically, not in the order they were processed in the C++ code. With such a confusion matrix written to a text file, when returning to the C++ world we can call the same \t{calculateAndLogMetrics} functions as for KNN and SVM, in the hope of equal treatment.
    
    \item \b{cnn\_classify\_board.py} is simple if the previous two scripts are understood. It just takes the board images and writes the predictions into a file.
\end{enumerate}

It was a good idea to include a CNN into the project since nowadays this is the go-to solution for image classification. The drawbacks were that I am far less familiar with Python, and every line of code required three googling. Even with such a powerful classifier we will see that the results are not the brightest, and this is due to the nature of the classification: 13 classes with small images and a tiny handmade dataset - no wonders should be expected.


% --------------------------------- File Handling --------------------------------- %

\subsection{File Handling}

The operations on file systems are often OS-specific and represent volatility. While some save/load operations can be done easily with higher-level functions, error handling requires complex logic. So it is worth grouping these operations into a separate module. The \t{file\_handling} submodule currently only contains one class, \t{FileHandler}.


% --- FileHandler --- %

\subsubsection{Class \t{FileHandler}}

Heavily relying on \t{Paths}, class \t{FileHandler} does operations on OpenCV images organized into a special hierarchy. We call folders\t{TEMPORARY\_FOLDER}, \t{TRAIN\_FOLDER}, \t{VALIDATION\_FOLDER} and \t{TEST\_FOLDER} label folders. This is because they all share the same internal structure: each one has 14 subfolders, containing images belonging to one known label. This is how images are prepared for classification. \t{TEMPORARY\_FOLDER} contains images that are labeled due to the structuring, but not yet associated with a data set.

The function \t{shuffleAndSplit} gathers all the images from all of the label folders in one vector, shuffles the vector and then distributes the images in training, validation and testing sets. The amount of images put into one ``partition'' depends on the input parameters, which should represent percentages, and are manipulable by the user. The moving operation is done optimally: instead of save and delete operations, the files are simply renamed. To avoid parsing (which often leads to hard-to-trace errors), I settled with some redundancy and used a structure called \t{PathInfo} to store all the information needed for the renaming operations ready to be used directly.

The function \t{clearAllImages} erases every image from \t{TEMPORARY\_FOLDER}, \t{TRAIN\_FOLDER}, \t{VALIDATION\_FOLDER}, \t{TEST\_FOLDER}, \t{BOARD\_FOLDER} and \t{GRAB\_FOLDER}. \t{readLabelFolderImages} is a high-level wrapper to get a list of image-label pairs from a labeled folder. \t{readBoardImages} is the same, just for the \t{BOARD\_FOLDER}, since this can be considered a ``simple folder'' (it does not obey the label folder structure). To ensure some uniformity throughout the project, \t{saveImage} provides an error-handling wrapper for OpenCV's \t{imwrite}, while \t{boardImageName} composes what name should a given cell image have based on its row and column in the main coordinate system.



% -------------------------------------------------------------------------------------------------- %
% --------------------------------- Validation and Response Module --------------------------------- %
% -------------------------------------------------------------------------------------------------- %

\section{Validation and Response Module}



All files of the VAR module can be found in folder \t{var}. The entry point is \t{ValidationAndResponse}. This is a \t{Q\_OBJECT}, since it has signals and slots to communicate with UA. Its three functionalities are implemented by its three slots, namely \t{validateMoveSlot}, \t{discardMoveSlot} and \t{newGameSlot}. When these are done with their job, they send replies through \t{validateMoveReplySignal}, \t{discardMoveReplySignal} and \t{newGameReplySignal}.

While the naming of the files in this module might not be the most intuitive, this is the only idea I had to split the functionalities without overcomplicating it. The idea is that we have a \t{Validator} structure, which contains the previous and current board (\t{prevBoard}, \t{currBoard}) and metadata. Metadata is information that needs to be kept track of to validate the moves or compose the long algebraic notation. This information is:

\begin{itemize}
    \item \b{turn}: the current color (whose turn it is)
    \item \b{moveCount}: incremented after each player moves once, could be part of the notation
    \item \b{enPassantCol}: the column of the board where en passant could be possible in the next move due to a pawn moving two blocks ahead. -1 if not possible
    \item \b{castle}: information about the eligibility of castling (did kings or rooks move).
\end{itemize}

It is assumed that the first board position (the value \t{prevBoard} will initially get) is the position where normal chess games start. This could be changed by changing field \t{initialSetup}. Besides logging, \t{Validatior} offers one method, which is \t{validateBoard}. This is where the magic happens.

We can only have one king of each color, this can be instantly checked, without any other prerequisites. The \t{Validator} instance is saved, so if the move is illegal, it can be set back to its original state, with the old \t{currBoard}, \t{prevBoard} and \t{metadata}. \t{Move.cpp}'s \t{processMove} is called, this validates the move using the information we have. If the move is valid, we can change turns and reset \t{enPassantCol} if it was not changed.

The function \t{processMove} gathers the changes we talked about in Subsection \ref{analysis:approach} and in Figure \ref{fig:move}. Zero changes mean no movement. We have to check for the status of both kings. This is done in \t{getKingSituation}. The final part of this code can be found in Appendix \ref{appendix:kingsituation}. If there are no attackers of the king, it might be free or in stalemate, otherwise it could be in check or checkmate. Extracting the king's situation is crucial because one cannot put himself/herself in check, and knowing if the king is in check is needed to evaluate a possible castleing move.

Based on the number of changes, the flow continues as shown in Figure \ref{fig:varflow}. In \t{\_isXMove} \t{X} could be: \t{WhiteMove}, \t{BlackMove}, \t{WhiteCapture}, \t{BlackCapture}, \t{WhiteEnPassant}, \t{BlackEnPassant}, \t{WhiteCastle}, \t{BlackCastle}. For better readability, changes are captured in a vector of \t{Change} structures. These structures contain the row and the column of the cell we are talking about, the previous piece's encoding (\t{prev}) and the same for the current piece (\t{curr}). There are also some inline functions defined for better readability, which are just wrappers around the encoding checking.

\begin{figure}[H]
    \centering
    \includegraphics[width=0.5\textwidth]{figs/varflow.png}
    \caption{Flow of execution in VAR}
    \label{fig:varflow}
\end{figure}

An example \t{\_isX} function's flow looks like:
\begin{enumerate}
    \item the turn is checked
    \item the setup shown in Figure \ref{fig:move} is ``forced''
    \item one or more \t{canPieceAttackCell} calls are made if necessary
    \item special checks are made
\end{enumerate}

Turn checking is simply checking if it is a white move that is white's turn, and vice versa. Forcing the setup refers to putting changes into ``slots'' \t{a} \t{b} etc. to have a one-to-one resemblence with the figure previously mentioned. This is important to further narrow down possibilities. Then we can check if indeed we can get the move listed as one of the possibilities. However, having the right pieces is not enough, the move has to be made according to the rules discussed in \ref{chessrules}. For this, we call \t{Attack}'s \t{canPieceAttackCell} on pieces that are involved.

The \t{canPieceAttackCell} function checks that we do not attack an allied piece, and delegates the checking to piece-specific functions. Pieces bishop and rook have their movement checked in a for loop. The moving path must be free, which is checked by \t{\_checkPath}. The queen is a combination of these two. The knight has a hardcoded pattern and can jump over other pieces. The pawn is a special case as well, since we have to check for simple forward move, double forward move, en passant left and right.

Going back to \t{Move}, special checks refer to cases which modify the metadata: movement of pawn may allow en passant next turn, movement of rook or king might disallow castling from now on, promotion has to be checked, and castling is hard-coded.

